% algebre2.tex --- exemples texecole, niveau collège
% !TeX program = lualatex
% Compilation : lualatex --shell-escape (les calculs exécutent Python)
\documentclass[marges=20, taille=11, precision=3]{texecole}
\begin{document}

soit h1 = <bleu marine gras section num>
soit h2 = <bleu italique sous-section>

<h1>Algèbre — collège

<h2>Fractions, pas à pas

<Calcule>1/2 + 1/3

<Calcule>7/6 - 3/4

<Calcule>3/4 * 2/5

<Calcule>1/2 : 3/4

<Simplifie>la fraction 84/60

<h2>Puissances de dix et notation scientifique

<Écris>le nombre 45 600 en notation scientifique

<Écris>le nombre 0,00032 en notation scientifique

<h2>Divisibilité et division euclidienne

<Vérifie>si 456 est divisible par 3

<Vérifie>si 4567 est divisible par 9

<Vérifie>si 1234 est divisible par 4

<Effectue>la division euclidienne de 47 par 5

<Calcule>le PGCD de 84 et 60

<h2>Programmes de calcul

<Applique>le programme de calcul à 5 {
	choisir un nombre
	ajouter 3
	multiplier par 2
	soustraire 4
}

<Applique>le programme de calcul à -2 {
	choisir un nombre
	élever au carré
	ajouter 1
}

<Exprime>le programme de calcul en fonction de x {
	choisir un nombre
	ajouter 3
	multiplier par 2
	soustraire 4
}

<h2>Équations du premier degré

<Résous>l'équation 3x + 5 = 17

<Résous>l'équation 2(x - 1) = x + 4

\end{document}
